\clearpage

% --------------------------------------------------------------------------------------------

\section {Summary}
\label{sec:summary}

In this technical note,
we have presented the results of a 3-flavour $\nu_{\mu}$-disappearance analysis,
performed by the {\em VALOR} group,
on the combined Run1+2+3+4 (6.5700$\times$10$^{20}$ POT) dataset.
This analysis is an update of the work presented in
Refs. \cite{OurRun1DisappAnalysis}, 
\cite{OurRun12DisappAnalysis}, 
\cite{OurRun12DisappAnalysis3Flavour},
\cite{OurRun123DisappAnalysis3Flavour} and 
\cite{OurRun1234aDisappAnalysis3Flavour}
with improved methods and an enlarged dataset.

This analysis predicts \NskNoOscExpectation $\pm$ \NskNoOscExpectationError (syst)
1 $\mu$-like ring events in SuperK in the absence of any oscillation,
but only \NskObserved were observed. The observed deficit has a strong energy dependence;
the ratio of observed to expected, under the no-oscillation hypothesis, is
$\sim$\NskNoOscRatioLow\% $<$ 0.5 GeV,
$\sim$\NskNoOscRatioMid\% between 0.5 and 1 GeV and
$\sim$\NskNoOscRatioHig\% $>$ 1 GeV.

A $\nu_{\mu}$-disappearance analysis was performed
in a framework of 3-flavour oscillations including matter effects in constant-density matter.
In this analysis, 41 neutrino flux, neutrino cross section,
efficiency, energy scale, initial-state nuclear environment and final-state rescattering,
and oscillation parameter systematics were taken into account. The oscillation parameters 
$\sin^{2}\theta_{12}$, $\sin^{2}\theta_{13}$, $\Delta m^{2}_{21}$ and $\delta_{CP}$ were 
also treated as nuisance parameters. 
The observed reconstructed energy spectrum of 1 $\mu$-like ring events was fitted,
and separate fits were made for the normal and the inverted mass hierarchies. In these fits,
$\sin^{2}\theta_{23}$ and either $|\Delta m^{2}_{32}|$ (normal hierarchy) or $|\Delta m^{2}_{31}|$ (inverted hierarchy)
were allowed to float.
The oscillation parameters $\sin^{2}\theta_{12}$, $\sin^{2}\theta_{13}$, $\Delta m^{2}_{21}$ and $\delta_{CP}$
and all 41 systematic parameters considered in this analysis were also allowed to float in the fits.

The 3-flavour $\nu_{\mu}$-disappearance fit of the combined Run 1+2+3+4 dataset,
assuming the normal mass hierarchy,
gives
\resultNormalDmsq and \resultNormalMixing
as the best-fit values (p-value = \resultNormalPvalue).
Assuming the inverted mass hierarchy, the fit yields
\resultInvertedDmsq and \resultInvertedMixing
as the best-fit values (p-value = \resultInvertedPvalue).

Confidence regions in the parameter spaces $\sin^{2}\theta_{23}$ - $|\Delta m^{2}_{32}|$ (normal hierarchy) and
$\sin^{2}\theta_{23}$ - $|\Delta m^{2}_{31}|$ (inverted hierarchy)
have been shown using both the constant-$\Delta\chi^2$ method and the method of Feldman and Cousins.
Using the new procedure described in Sec.\ \ref{sec:contour_construction}, the 68\% (90\%) confidence intervals 
in $\sin^{2}\theta_{23}$ and either $|\Delta m^{2}_{32}|$ (normal hierarchy) or $|\Delta m^{2}_{31}|$ (inverted hierarchy) 
can be summarized as
% This is 1D FC
\resultNormalMixingOneSigmaFCLo\xspace(\resultNormalMixingTwoSigmaFCLo)
$< \sin^{2}\theta_{23} <$
\resultNormalMixingOneSigmaFCHi\xspace(\resultNormalMixingTwoSigmaFCHi)
and
\resultNormalDmsqOneSigmaFCLo$\times$10$^{-3}$ (\resultNormalDmsqTwoSigmaFCLo$\times$10$^{-3}$) eV$^{2}$/c$^{4}$
$< |\Delta m^{2}_{32}| < $
\resultNormalDmsqOneSigmaFCHi$\times$10$^{-3}$ (\resultNormalDmsqTwoSigmaFCHi$\times$10$^{-3}$) eV$^{2}$/c$^{4}$
for the normal mass hierarchy,
and
\resultInvertedMixingOneSigmaFCLo\xspace(\resultInvertedMixingTwoSigmaFCLo)
$< \sin^{2}\theta_{23} <$
\resultInvertedMixingOneSigmaFCHi\xspace(\resultInvertedMixingTwoSigmaFCHi)
and
\resultInvertedDmsqOneSigmaFCLo$\times$10$^{-3}$ (\resultInvertedDmsqTwoSigmaFCLo$\times$10$^{-3}$) eV$^{2}$/c$^{4}$
$< |\Delta m^{2}_{31}| < $
\resultInvertedDmsqOneSigmaFCHi$\times$10$^{-3}$ (\resultInvertedDmsqTwoSigmaFCHi$\times$10$^{-3}$) eV$^{2}$/c$^{4}$
for the inverted mass hierarchy.
% % This is const dchisq
% \resultNormalMixingOneSigmaLo\xspace(\resultNormalMixingTwoSigmaLo)
% $< \sin^{2}\theta_{23} <$
% \resultNormalMixingOneSigmaHi\xspace(\resultNormalMixingTwoSigmaHi)
% and
% \resultNormalDmsqOneSigmaLo$\times$10$^{-3}$ (\resultNormalDmsqTwoSigmaLo$\times$10$^{-3}$) eV$^{2}$/c$^{4}$
% $< |\Delta m^{2}_{32}| < $
% \resultNormalDmsqOneSigmaHi$\times$10$^{-3}$ (\resultNormalDmsqTwoSigmaHi$\times$10$^{-3}$) eV$^{2}$/c$^{4}$
% for the normal mass hierarchy,
% and
% \resultInvertedMixingOneSigmaLo\xspace(\resultInvertedMixingTwoSigmaLo)
% $< \sin^{2}\theta_{23} <$
% \resultInvertedMixingOneSigmaHi\xspace(\resultInvertedMixingTwoSigmaHi)
% and
% \resultInvertedDmsqOneSigmaLo$\times$10$^{-3}$ (\resultInvertedDmsqTwoSigmaLo$\times$10$^{-3}$) eV$^{2}$/c$^{4}$
% $< |\Delta m^{2}_{31}| < $
% \resultInvertedDmsqOneSigmaHi$\times$10$^{-3}$ (\resultInvertedDmsqTwoSigmaHi$\times$10$^{-3}$) eV$^{2}$/c$^{4}$
% for the inverted mass hierarchy.

%Numerous comparisons have been made of the best-fit values and spectra 
%for different running periods.
Comparisons have been presented of the confidence regions obtained using different construction methods.
%and from fits to different datasets.
Finally, comparisons of our contours with
recent measurements from MINOS and SuperK
and the expected T2K sensitivity 
for 6.570$\times$10$^{20}$ POT
have also been made. 

